\documentclass[compress]{beamer}
\usetheme{Warsaw}
\usepackage{fancybox}
\usepackage[pstricks]{bclogo}
\usepackage{tcolorbox}
\usepackage{listings}

\usepackage{fontspec}
\setsansfont{Acathist}

\usecolortheme[RGB={200,0,0}]{structure}
\setbeamertemplate{navigation symbols}{}
\setbeamertemplate{headline}{}

\definecolor{myblue}{RGB}{46, 44, 139}
\definecolor{myred}{RGB}{180, 40, 40}
\definecolor{myviolet}{RGB}{110, 40, 140}
\definecolor{kw}{RGB}{0,0,200}
\definecolor{cm}{RGB}{0,128,0}
\definecolor{ty}{RGB}{43,145,175}


\tcbset{
	noparskip,
	coltext=black,
	coltitle=white,
	theorem/.style={
		colback=myred!9,
		colframe=myred,
	},
	example/.style={
		colback=green!9,
		colframe=green!40!black,
	},
	definition/.style={
		colback=myblue!9,
		colframe=myblue,
	},
}



\title{THERE IS NO LARGEST PRIME NUMBERS}
\author[]{IRIMANOA Nampoina Hiarijaona}
\institute{CRIAT MAdagascar \\ Center Antananarivo}
\date{\today}

\begin{document}
\begin{frame}[plain]
    \maketitle\end{frame}

\begin{frame}[plain]
	\frametitle{Outlines}
	\tableofcontents
\end{frame}

\section{Prime numbers}
\subsection{What are prime numbers ?}
\begin{frame}{What Are Prime Number ?}
	\begin{tcolorbox}[definition, title=Definition]
		A prime number is a number that has exactly two divisors
	\end{tcolorbox}
	\begin{tcolorbox}[example, title=Exemple]
		\begin{itemize}
			\item 2 is prime (two divisors: 1 and 2)
			\item 3 is prime (two divisors: 1 and 3)
			\item 4 is not prime (three divisors: 1, 2 and 4)
		\end{itemize}
	\end{tcolorbox}
\end{frame}

\subsection{There is no largest prime numbers}
\begin{frame}{There Is No Largest Prime Numbers}
	\transdissolve
	\framesubtitle{The proof uses \textit{redictio ad absurdum}.}
	
	\begin{tcolorbox}[theorem, title=Theorem]
		There is no largest prime numbers
	\end{tcolorbox}
	
	\begin{bclogo}[couleur=myred!9,couleurBarre=myred,arrondi=0.2,logo=\bclampe,noborder=true,marge=10]{  Proof}
		\begin{enumerate}
			\item Suppose $p$ is the largest prime number.
			\item Let $q$ the product of the first $p$ numbers.
			\item Then $q + 1$ is not divisilbe be any of them.
			\item But $q + 1$ is greater than $1$, thus divisible by some prime numbers not in the the first $p$ prime number.\qedhere 
		\end{enumerate}
	\end{bclogo}
\end{frame}

\begin{frame}[fragile]
	\frametitle{An Algorithm For Finding Prime Numbers.}
	\begin{verbatim}
int main (void)
{
  std::vector<bool> is_prime (100, true);
  for (int i = 2; i < 100; i++)
  if (is_prime[i])
  {
    std::cout << i << " ";
    for (int j = i; j < 100; is_prime [j] = false, j+=i);
  }
  return 0;
}
    \end{verbatim}
		Note the use of \alert{std::} .
\end{frame}

\end{document}
